\section{Discussion}
\label{sec:discussion}

\paragraph{What the band contributes.}
Limiting a ranker to a $10\mm$ native ball removes evidence that a wider view can recover. Restoring one thin band recovers about seven tenths of the CPM difference to the full field in two training runs, and the effect carries over to a single-stream 3D network. Sparse access, however it is read, recovers much less. Showing the image network the exact voxels behind the 384 samples closes less than half of the gap, and attention over sample directions adds little to a late readout. Reading a few well-chosen places, even with their geometry, is therefore not a substitute for seeing a contiguous region. A useful follow-up would vary the arrangement of visible voxels at a fixed count, separating coverage from continuity.

\paragraph{What ``exterior'' means for detection.}
The band lies outside a ball centred on the \emph{proposal}, not on the lesion. Large lesions and off-centre proposals put lesion tissue inside it, and this is where the largest gains occur, using a grouping fixed before scoring. Small benefits beyond the annotated lesion remain compatible with the data, especially at higher FP rates. The FROC changes could also come from better rejection of false proposals, even if a true proposal's own appearance were unchanged. Attributing gains to vessels, attachments or lesion margins would require other interventions or anatomical labels. For detector design, the result suggests checking how often proposals are displaced from the lesion centre before investing in long-range context modules.

\paragraph{Lesion size, proposal displacement and context.}
A $10\mm$ ball is wide enough to contain most small nodules, so for lesions below $6\mm$ the band adds only surroundings; all of them are in the never-reached group, whose gain at 1~FP/scan is small. For lesions of at least $10\mm$, the ball cuts through the nodule, and two thirds of them have lesion tissue in the band. Large lesions also attract more proposals: the 72 reached lesions have 221 matching proposals, whereas each never-reached lesion has exactly one. Many of these proposals lie away from the lesion centre, which is where the band restores lesion tissue. The context advantage reported for larger patches may therefore partly be a \emph{coverage} advantage for large lesions. Analyses of context in detection should report effects by lesion size, and patch sizes may be better chosen from the size distribution of the target lesions than from the scale of surrounding anatomy.

\paragraph{Limitations.}
Masking creates artificial air boundaries, and each model can adapt to them; the native-source and band conditions differ both in the voxels they show and in how these boundaries look after resampling. We tested one band, chosen to enclose the samples' support; other radii and equal-volume regions remain open. The sparse heads use unoriented axis pairs at one fusion point, so their results are specific to this design. In fit~2, three attention heads did not move from their initialization, which weakens the geometry evidence to fit~1 and the width-48 contrast. The native-source and single-stream comparisons each use one training seed, and the patient intervals do not capture variation across training runs. All fits share one candidate generator and one benchmark that has been studied before; an independent test would need an untouched full-scan cohort with a fixed proposal stream.

\section{Conclusion}
\label{sec:conclusion}

We asked what a wider view adds when a lung-nodule ranker scores fixed candidates, and answered it by changing only which native CT voxels the ranker can see. Restoring a thin, physically defined band around a $10\mm$ ball raises two-scale 3D CPM by about four points in two training runs and closes about 70\% of the gap to the full field; a second 3D architecture also gains. The dense band beats the exact voxels behind sparse samples in the same encoder, and adding sample geometry to sparse readouts does not close this gap. The largest gains occur where the band reaches annotated lesion tissue. Visibility interventions of this kind give a simple, architecture-independent way to ask which evidence a medical detector actually needs.
